\begingroup
\setlength{\LTcapwidth}{\textwidth}
\small
\captionsetup{font=normalsize}
\setlength{\tabcolsep}{4pt}
\begin{longtable}{@{}>{\raggedright\arraybackslash}p{1.750cm}>{\raggedright\arraybackslash}p{5.350cm}>{\raggedright\arraybackslash}p{6.338cm}@{}}
\caption{Pemetaan kebutuhan sistem terhadap bagian rancangan yang merealisasikannya.}
\label{tab:tab4.2_pemetaan_kebutuhan_rancangan} \\
\toprule
\textbf{Kebutuhan} & \textbf{Inti kebutuhan} & \textbf{Bagian rancangan} \\
\midrule
\endfirsthead

\multicolumn{3}{@{}l}{\small\emph{Tabel \thetable{} (lanjutan)}} \\
\toprule
\textbf{Kebutuhan} & \textbf{Inti kebutuhan} & \textbf{Bagian rancangan} \\
\midrule
\endhead

\bottomrule
\endlastfoot

KF-01 & Menerima catatan \emph{logbook} terstruktur & Subbab~\ref{subsec:bab4-input-output-modul-prediksi} dan Subbab~\ref{subsec:bab4-prapemrosesan-kualitas-data} \\

KF-02 & Menghitung besaran turunan berbasis fisiologi & Subbab~\ref{subsec:bab4-rekayasa-fitur-prediksi} \\

KF-03 & Memprediksi glukosa pada dua horizon beserta selang ketidakpastian & Subbab~\ref{subsec:bab4-pembentukan-sekuens-horizon}, Subbab~\ref{subsec:bab4-target-prediksi-rekonstruksi}, dan Subbab~\ref{subsec:bab4-ketidakpastian-prediksi} \\

KF-04 & Menyusun kueri yang dikondisikan pada glukosa terprediksi & Subbab~\ref{subsec:bab4-query-transformation-prediksi} dan Subbab~\ref{subsec:bab4-perbandingan-rag} \\

KF-05 & Menelusuri korpus pedoman dan mengembalikan potongan ber-penunjuk halaman & Subbab~\ref{subsec:bab4-kurasi-basis-pengetahuan}, Subbab~\ref{subsec:bab4-segmentasi-dokumen-metadata}, dan Subbab~\ref{subsec:bab4-desain-penelusuran} \\

KF-06 & Menyusun rekomendasi berbahasa Indonesia yang bersandar pada dokumen & Subbab~\ref{subsec:bab4-grounded-generation} \\

KF-07 & Menyajikan kondisi klinis terstruktur beserta peringatan dini dan divergensi & Subbab~\ref{subsec:bab4-representasi-kondisi-klinis} dan Subbab~\ref{subsec:bab4-struktur-assessment-output} \\

KF-08 & Mendukung alur termediasi dokter tanpa tindakan otonom & Subbab~\ref{subsec:bab4-guardrail-keselamatan} dan Subbab~\ref{subsec:bab4-konteks-penggunaan-doctor-mediated} \\

\midrule

KNF-01 & Keterterapan tanpa akselerator grafis & Subbab~\ref{subsec:bab4-arsitektur-sistem} dan Subbab~\ref{subsec:bab4-desain-penelusuran} \\

KNF-02 & Keterjelasan prediksi dan rekomendasi & Subbab~\ref{subsec:bab4-ketidakpastian-prediksi} dan Subbab~\ref{subsec:bab4-keterlacakan-verifikasi-sitasi} \\

KNF-03 & Keamanan konten keluaran & Subbab~\ref{subsec:bab4-guardrail-keselamatan} \\

KNF-04 & Ketahanan terhadap data dan evidence yang tidak memadai & Subbab~\ref{subsec:bab4-fallback-retrieval-generation} dan Subbab~\ref{subsec:bab4-integrasi-fallback-status-data} \\

KNF-05 & Keluaran berbahasa Indonesia & Subbab~\ref{subsec:bab4-grounded-generation} \\

KNF-06 & Keamanan klinis dan penentu akhir pada dokter & Subbab~\ref{subsec:bab4-guardrail-keselamatan} dan Subbab~\ref{subsec:bab4-konteks-penggunaan-doctor-mediated} \\

KNF-07 & Keterpeliharaan melalui pemisahan komponen & Subbab~\ref{subsec:bab4-arsitektur-sistem} \\

KNF-08 & Keterlacakan sitasi sampai halaman sumber & Subbab~\ref{subsec:bab4-segmentasi-dokumen-metadata} dan Subbab~\ref{subsec:bab4-keterlacakan-verifikasi-sitasi} \\

KNF-09 & Biaya rendah dan keterpindahan lingkungan & Subbab~\ref{subsec:bab4-arsitektur-sistem} \\

KNF-10 & Waktu tanggap yang memungkinkan penggunaan saat konsultasi & Subbab~\ref{subsec:bab4-aliran-data-end-to-end} \\
\end{longtable}
\endgroup
